\small
\vspace{0.5cm}

1.- Dado el sistema no lineal SISO

\begin{equation*}
    \begin{aligned}
        \dot{x}&= Ax+f(x,t)+Bu \\
        y&= Cx
    \end{aligned}
\end{equation*}

donde

\begin{equation*}
    A = \begin{bmatrix}
        0 & 1 & 0 & \cdots & 0 \\
        0 & 0 & 1 & \cdots & 0 \\
        \vdots & \vdots & \vdots & \ddots & \vdots \\
        0 & 0 & 0 & \cdots & 1 \\
        0 & 0 & 0 & \cdots & 0
    \end{bmatrix}, \qquad
    B = \begin{bmatrix}
        0 \\
        0 \\
        \vdots \\
                0 \\
        1
    \end{bmatrix}, \qquad
    C = \begin{bmatrix}
        1 & 0  & \cdots & 0 & 0
    \end{bmatrix}, \qquad
    f(x,t) = \begin{bmatrix}
        f_1(x,t) \\
        f_2(x,t) \\
        \vdots \\
        f_{n-1}(x,t) \\
        f_n(x,t)
    \end{bmatrix}
\end{equation*}

Usando el método de elipsoides atractivas:

\begin{itemize}
    \item Verifique que la ley de control
    
    \begin{equation*}
        u = -Kx 
    \end{equation*}

    estabiliza el sistema en el sentido de Lyapunov.

    \textbf{Solución:}

Definimos la matriz de lazo cerrado como
\begin{equation*}
    A_{lc} = A - BK
\end{equation*}
por lo que entonces el sistema es
\begin{equation*}
    \dot{x} = A_{lc}x + f(x,t)
\end{equation*}
y proponemos que la no linealidad está acotada por
\begin{equation*}
    \|f(x,t)\| \leq D, \quad D>0
\end{equation*}

Proponemos la función de Lyapunov
\begin{equation*}
    V(x) = x^TPx
\end{equation*}
derivando
\begin{equation*}
    \dot{V}(x) = x^T(A_{lc}^TP + PA_{lc})x + 2x^TPf(x,t)
\end{equation*}

Ocupando la desigualdad de Young
\begin{equation*}
    2ab \leq \epsilon a^2 + \frac{1}{\epsilon}b^2
\end{equation*}
donde \(a = \|Px\|\) y \(b = \|f(x,t)\|\), se tiene que
\begin{equation*}
    2x^TPf(x,t) \leq \epsilon x^TP^2x + \frac{1}{\epsilon}\|f(x,t)\|^2
\end{equation*}

Sustituyendo la cota propuesta para la no linealidad, obtenemos
\begin{equation*}
    2x^TPf(x,t) \leq \epsilon x^TP^2x + \frac{1}{\epsilon}D^2   
\end{equation*}

Sustituyendo en la derivada de la función de Lyapunov, obtenemos
\begin{equation*}
    \dot{V}(x) \leq x^T(A_{lc}^TP + PA_{lc} + \epsilon P^2)x + \frac{1}{\epsilon}D^2
\end{equation*}

Queremos que
\begin{equation*}
    \dot{V} \leq - \alpha V+\frac{1}{\epsilon}D^2 = -\alpha x^TPx + \frac{1}{\epsilon}D^2
\end{equation*}

Se propone entonces que
\begin{equation*}
x^{T}\left(A_{lc}^{T}P + PA_{lc} + \epsilon P^{2}\right)x \leq -\alpha x^{T}Px
\end{equation*}
\begin{equation*}
   x^T \left(A_{lc}^TP+PA_{lc}+\epsilon P^2 +\alpha P\right) x \leq 0
\end{equation*}

La expresión anterior no es lineal en \((P,K)\) porque \(A_{lc}=A-BK\) produce el producto \(PBK\), y el término \(\epsilon P^2\) también es no lineal. Se aplican los siguientes pasos para transformarla en una LMI. Se sustituyendo \(A_{lc} = A - BK\) en la condición:

\begin{equation*}
    (A - BK)^T P + P(A - BK) + \epsilon P^2 + \alpha P \leq 0
\end{equation*}
Desarrollando:
\begin{equation*}
    A^T P + P A - K^T B^T P - P B K + \epsilon P^2 + \alpha P \leq 0
\end{equation*}

El término \(P B K\) es bilineal, para eliminarlo se propone el cambio de variable:
\begin{equation*}
    Y := K P \quad \Longleftrightarrow \quad K = Y P^{-1}
\end{equation*}
Sustituyendo en el término bilineal:
\begin{equation*}
    P B K = P B (Y P^{-1}) = P B Y P^{-1}
\end{equation*}
y su transpuesto:
\begin{equation*}
    K^T B^T P = (Y P^{-1})^T B^T P = P^{-T} Y^T B^T P
\end{equation*}

Se multiplica la desigualdad por \(P^{-1}\) a izquierda y a derecha. Definiendo \(Q = P^{-1}\), se obtiene:
\begin{equation*}
    Q A^T + A Q - Q (P B Y P^{-1}) Q - Q (P^{-1} Y^T B^T P) Q + \epsilon Q P^2 Q + \alpha Q \leq 0
\end{equation*}
Simplificando cada término:
\begin{equation*}
    Q (P B Y P^{-1}) Q = (Q P) B Y (P^{-1} Q) = I \cdot B Y \cdot I = B Y
\end{equation*}
\begin{equation*}
    Q (P^{-1} Y^T B^T P) Q = (Q P^{-1}) Y^T B^T (P Q) = I \cdot Y^T B^T \cdot I = Y^T B^T
\end{equation*}
\begin{equation*}
    Q P^2 Q = Q P \cdot P Q = I \cdot I = I
\end{equation*}
Por lo tanto, la condición se transforma en:
\begin{equation*}
    A Q + Q A^T - B Y - Y^T B^T + \epsilon I + \alpha Q \leq 0
\end{equation*}
donde \(Q = P^{-1} \succ 0\). 

La desigualdad anterior contiene el término \(\epsilon I\) que no depende de las variables. Para incorporarlo a una LMI, se aplica el complemento de Schur. Para una matriz simétrica particionada
\begin{equation*}
    M = \begin{bmatrix} M_{11} & M_{12} \\ M_{12}^T & M_{22} \end{bmatrix},
\end{equation*}
se tiene que \(M \prec 0\) si y solo si \(M_{22} \prec 0\) y \(M_{11} - M_{12} M_{22}^{-1} M_{12}^T \prec 0\). Considerando que la desigualdad
\begin{equation*}
    A Q + Q A^T - B Y - Y^T B^T + \alpha Q + \epsilon I \leq 0
\end{equation*}
se puede reescribir como:
\begin{equation*}
    A Q + Q A^T - B Y - Y^T B^T + \alpha Q + \frac{1}{\delta} (\delta I)^T (\epsilon I)^{-1} (\delta I) \leq 0
\end{equation*}
donde \(\delta > 0\) es un parámetro de escalamiento. Aplicando el complemento de Schur con:
\begin{equation*}
    M_{11} = A Q + Q A^T - B Y - Y^T B^T + \alpha Q + \epsilon I, \quad
    M_{12} = \delta I, \quad
    M_{22} = -\epsilon I,
\end{equation*}
se obtiene la LMI:
\begin{equation*}
    \begin{bmatrix}
        A Q + Q A^T - B Y - Y^T B^T + \alpha Q + \epsilon I & \delta I \\
        \delta I & -\epsilon I
    \end{bmatrix} \leq 0.
\end{equation*}

Para mantener la simetría y la estructura estándar de la LMI en la variable \(P\), se puede aplicar el complemento de Schur directamente sobre la desigualdad en \(P\). Partiendo de:
\begin{equation*}
    A^T P + P A - K^T B^T P - P B K + \epsilon P^2 + \alpha P \leq 0,
\end{equation*}
y usando \(Y = KP\), se tiene \(P B K = B Y\) y \(K^T B^T P = Y^T B^T\). Entonces:
\begin{equation*}
    A P + P A^T - B Y - Y^T B^T + \epsilon P^2 + \alpha P \leq 0.
\end{equation*}
El término \(\epsilon P^2\) se puede linealizar aplicando el complemento de Schur con un parámetro \(\delta > 0\):
\begin{equation*}
    \epsilon P^2 = (\delta P)^T (\epsilon I)^{-1} (\delta P) \quad \text{si } \delta = \sqrt{\epsilon}.
\end{equation*}
Aplicando Schur a:
\begin{equation*}
    M_{11} = A P + P A^T - B Y - Y^T B^T + \alpha P + \epsilon I, \quad
    M_{12} = \delta P, \quad
    M_{22} = -\epsilon I,
\end{equation*}
se obtiene finalmente:
\begin{equation*}
    \begin{bmatrix}
        A P + P A^T + B Y + Y^T B^T + \alpha P + \epsilon I & \delta P \\
        \delta P & -\epsilon I
    \end{bmatrix} \leq 0,
\end{equation*}
sujeta a
\begin{equation*}
    P \succ 0, \qquad \alpha > 0, \qquad \epsilon > 0, \qquad \delta > 0.
\end{equation*}
Una vez resuelta la LMI y obtenidos \((P, Y)\), la ganancia se recupera como
\begin{equation*}
    K = Y (P)^{-1}.
\end{equation*}

Para verificar la elipsoide atractiva de
\begin{equation*}
    \dot{V} \leq -\alpha V + \frac{1}{\epsilon}D^2
\end{equation*}
multiplicando por \(e^{\alpha t}\)
\begin{equation*}
    e^{\alpha t}\dot{V} + \alpha e^{\alpha t}V \leq \frac{1}{\epsilon}D^2e^{\alpha t}
\end{equation*}
\begin{equation*}
    \frac{d}{dt}(e^{\alpha t}V) \leq \frac{1}{\epsilon}D^2e^{\alpha t}
\end{equation*}
Integrando de \(0\) a \(t\)
\begin{equation*}
    e^{\alpha t}V(t) - V(0) \leq \frac{D^2}{\epsilon \alpha} (e^{\alpha t}-1)
\end{equation*}
\begin{equation*}
    V(t) \leq V(0)e^{-\alpha t} + \frac{D^2}{\epsilon \alpha}(1-e^{-\alpha t})
\end{equation*}
Considerando el límite cuando \(t \to \infty\)
\begin{equation*}
    V_{superior} \leq \frac{D^2}{\epsilon \alpha}
\end{equation*}

Ahora, considerando la desigualdad de Rayleigh, para \(P\) simétrica y definida positiva, y para todo \(x\), se tiene que
\begin{equation*}
    \lambda_{\min}(P)\|x\|^2 \leq x^TPx \leq \lambda_{\max}(P)\|x\|^2
\end{equation*}
Considerando que \(A=P\)
\begin{equation*}
    \lambda_{\min}(P)\|x\|^2 \leq x^TPx = V(x) \to \|x\| \leq \sqrt{\frac{V(x)}{\lambda_{\min}(P)}}
\end{equation*}
Combinando con el límite superior de \(V\), obtenemos
\begin{equation*}
    \|x\| \leq \sqrt{\frac{D^2}{\epsilon \alpha \lambda_{\min}(P)}}
\end{equation*}
por lo que la elipsoide atractiva es
\begin{equation*}
    \mathcal{E}=\left\{x \in \mathbb{R}^n : x^TPx \leq \frac{D^2}{\epsilon \alpha} \right\}
\end{equation*}



    \item Proponga un controlador por realimentación de estados que permita realizar el seguimiento de trayectorias por parte de la salida
    
    Sea el sistema no lineal SISO:
\begin{equation*}
    \begin{aligned}
        \dot{x} &= A x + f(x,t) + B u \\
        y &= C x
    \end{aligned}
\end{equation*}
donde \(x \in \mathbb{R}^n\), \(u \in \mathbb{R}\), \(y \in \mathbb{R}\), \(A \in \mathbb{R}^{n \times n}\), \(B \in \mathbb{R}^{n \times 1}\), \(C \in \mathbb{R}^{1 \times n}\), y \(f(x,t)\) es una no linealidad acotada.

Se desea que la salida \(y\) siga una referencia \(r(t)\). Se define el error de seguimiento de la salida como:
\begin{equation*}
    e = r - C x.
\end{equation*}
Para garantizar convergencia asintótica del error, se agrega el estado integral:
\begin{equation*}
    x_{n+1} = \int e \, dt = \int (r - C x) \, dt \quad \Rightarrow \quad \dot{x}_{n+1} = e = r - C x.
\end{equation*}

Para el caso general, los errores sucesivos se definen como las derivadas sucesivas del error de salida:
\begin{equation*}
    e_i = \frac{d^{i-1}}{dt^{i-1}}(r - y), \quad i = 1, 2, \dots, n.
\end{equation*}
Esta definición es válida cuando el sistema tiene grado relativo \(n\), es decir, cuando la entrada \(u\) aparece por primera vez en la \(n\)-ésima derivada de la salida.

En el caso particular en que la planta está en forma canónica controlable  y \(C = \begin{bmatrix} 1 & 0 & \cdots & 0 \end{bmatrix}\), se tiene \(y = x_1\), y los errores de seguimiento se reducen a:
\begin{equation*}
    e_1 = r - x_1, \quad e_2 = \dot{r} - x_2, \quad \dots, \quad e_n = r^{(n-1)} - x_n,
\end{equation*}
donde \(r^{(i-1)}\) es la derivada de orden \(i-1\) de la referencia.

Se define el vector de errores aumentado:
\begin{equation*}
    e = \begin{bmatrix} e_1 \\ e_2 \\ \vdots \\ e_n \\ x_{n+1} \end{bmatrix} \in \mathbb{R}^{n+1}.
\end{equation*}

Se propone la ley de control por realimentación de estados sobre el error:
\begin{equation*}
    u = -K e, \quad K = \begin{bmatrix} k_1 & k_2 & \cdots & k_n & k_i \end{bmatrix} \in \mathbb{R}^{1 \times (n+1)}.
\end{equation*}

Derivando los errores y sustituyendo la dinámica del sistema, se obtiene la dinámica del error en forma matricial:
\begin{equation*}
    \dot{e} = A_e e + B_e u + D_e f(x,t),
\end{equation*}
donde las matrices \(A_e\), \(B_e\) y \(D_e\) se construyen a partir de las matrices originales del sistema. Para el caso de una planta en forma canónica controlable con un integrador adicional, se tiene:
\begin{equation*}
    A_e = \begin{bmatrix}
        0 & 1 & 0 & \cdots & 0 & 0 \\
        0 & 0 & 1 & \cdots & 0 & 0 \\
        \vdots & \vdots & \vdots & \ddots & \vdots & \vdots \\
        0 & 0 & 0 & \cdots & 1 & 0 \\
        0 & 0 & 0 & \cdots & 0 & 0 \\
        1 & 0 & 0 & \cdots & 0 & 0
    \end{bmatrix} \in \mathbb{R}^{(n+1) \times (n+1)}, \quad
    B_e = \begin{bmatrix} 0 \\ 0 \\ \vdots \\ 0 \\ 1 \\ 0 \end{bmatrix} \in \mathbb{R}^{(n+1) \times 1},
\end{equation*}
\begin{equation*}
    D_e = \begin{bmatrix}
        I_{n \times n} \\
        0_{1 \times n}
    \end{bmatrix} \in \mathbb{R}^{(n+1) \times n}.
\end{equation*}
La no linealidad \(f(x,t)\) actúa como una perturbación \textit{mismatched}, ya que \(D_e \neq B_e\), lo que impide su cancelación directa mediante la entrada \(u\). Sustituyendo la ley de control \(u = -K e\), el sistema en lazo cerrado es:

\begin{equation*}
    \dot{e} = A_{lc} e + D_e f(x,t), \quad A_{lc} = A_e - B_e K.
\end{equation*}

Se supone que la perturbación está acotada por:
\begin{equation*}
    \|f(x,t)\| \leq D, \quad D > 0.
\end{equation*}

Se propone la función candidata de Lyapunov:
\begin{equation*}
    V(e) = e^T P e, \quad P > 0.
\end{equation*}

Derivando:
\begin{equation*}
    \dot{V}(e) = e^T(A_{lc}^T P + P A_{lc})e + 2e^T P D_e f(x,t).
\end{equation*}

Ocupando la desigualdad de Young:
\begin{equation*}
    2ab \leq \epsilon a^2 + \frac{1}{\epsilon}b^2
\end{equation*}
donde \(a = \|D_e^T P e\|\) y \(b = \|f(x,t)\|\), se tiene que:
\begin{equation*}
    2e^T P D_e f(x,t) \leq \epsilon e^T P D_e D_e^T P e + \frac{1}{\epsilon}\|f(x,t)\|^2.
\end{equation*}

Sustituyendo la cota de la no linealidad:
\begin{equation*}
    2e^T P D_e f(x,t) \leq \epsilon e^T P D_e D_e^T P e + \frac{1}{\epsilon}D^2.
\end{equation*}

Sustituyendo en la derivada de \(V\):
\begin{equation*}
    \dot{V}(e) \leq e^T\left( A_{lc}^T P + P A_{lc} + \epsilon P D_e D_e^T P \right)e + \frac{1}{\epsilon}D^2.
\end{equation*}

Se desea que:
\begin{equation*}
    \dot{V}(e) \leq -\alpha V(e) + \frac{1}{\epsilon}D^2 = -\alpha e^T P e + \frac{1}{\epsilon}D^2.
\end{equation*}

Por lo tanto, se requiere:
\begin{equation*}
    A_{lc}^T P + P A_{lc} + \epsilon P D_e D_e^T P + \alpha P \leq 0.
\end{equation*}

Se aplican los siguientes pasos para transformar la expresion de una BMI en una LMI. Sustituyendo \(A_{lc} = A_e - B_e K\):
\begin{equation*}
    (A_e - B_e K)^T P + P(A_e - B_e K) + \epsilon P D_e D_e^T P + \alpha P \leq 0.
\end{equation*}
Desarrollando:
\begin{equation*}
    A_e^T P + P A_e - K^T B_e^T P - P B_e K + \epsilon P D_e D_e^T P + \alpha P \leq 0.
\end{equation*}

Se propone el cambio de variable:
\begin{equation*}
    Y := K P \quad \Longleftrightarrow \quad K = Y P^{-1}.
\end{equation*}
Sustituyendo \(K = Y P^{-1}\) en el término bilineal:
\begin{equation*}
    P B_e K = P B_e (Y P^{-1}) = P B_e Y P^{-1}.
\end{equation*}

Definiendo \(Q := P^{-1}\), se obtiene:

\begin{equation*}
    Q A_e^T + A_e Q - Q (P B_e Y P^{-1}) Q - Q (P^{-1} Y^T B_e^T P) Q + \epsilon Q P D_e D_e^T P Q + \alpha Q \leq 0.
\end{equation*}
Simplificando cada término usando \(Q = P^{-1}\):
\begin{equation*}
    Q (P B_e Y P^{-1}) Q = (Q P) B_e Y (P^{-1} Q) = I \cdot B_e Y \cdot I = B_e Y,
\end{equation*}
\begin{equation*}
    Q (P^{-1} Y^T B_e^T P) Q = (Q P^{-1}) Y^T B_e^T (P Q) = I \cdot Y^T B_e^T \cdot I = Y^T B_e^T,
\end{equation*}
\begin{equation*}
    Q P D_e D_e^T P Q = (Q P) D_e D_e^T (P Q) = I \cdot D_e D_e^T \cdot I = D_e D_e^T.
\end{equation*}
Por lo tanto, la condición se transforma en:
\begin{equation*}
    A_e Q + Q A_e^T - B_e Y - Y^T B_e^T + \epsilon D_e D_e^T + \alpha Q \leq 0,
\end{equation*}
donde \(Q = P^{-1} \succ 0\). 

Aplicando el complemento de Schur como:

\begin{equation*}
    M = \begin{bmatrix} M_{11} & M_{12} \\ M_{12}^T & M_{22} \end{bmatrix},
\end{equation*}
se tiene que \(M \prec 0\) si y solo si \(M_{22} \prec 0\) y \(M_{11} - M_{12} M_{22}^{-1} M_{12}^T \prec 0\). La desigualdad

\begin{equation*}
    A_e Q + Q A_e^T - B_e Y - Y^T B_e^T + \alpha Q + \epsilon D_e D_e^T \leq 0
\end{equation*}
se puede reescribir como:
\begin{equation*}
    A_e Q + Q A_e^T - B_e Y - Y^T B_e^T + \alpha Q + (\sqrt{\epsilon} D_e) (\sqrt{\epsilon} D_e)^T \leq 0.
\end{equation*}
Aplicando el complemento de Schur con:
\begin{equation*}
    M_{11} = A_e Q + Q A_e^T - B_e Y - Y^T B_e^T + \alpha Q, \quad
    M_{12} = \sqrt{\epsilon} D_e, \quad
    M_{22} = -I,
\end{equation*}
se obtiene la LMI:
\begin{equation*}
    \begin{bmatrix}
        A_e Q + Q A_e^T - B_e Y - Y^T B_e^T + \alpha Q & \sqrt{\epsilon} D_e \\
        \sqrt{\epsilon} D_e^T & -I
    \end{bmatrix} \leq 0.
\end{equation*}

Equivalentemente, definiendo \(\delta = \sqrt{\epsilon}\), la LMI se escribe como:
\begin{equation*}
    \begin{bmatrix}
        A_e Q + Q A_e^T - B_e Y - Y^T B_e^T + \alpha Q & \delta D_e \\
        \delta D_e^T & -I
    \end{bmatrix} \leq 0,
\end{equation*}
sujeta a:
\begin{equation*}
    Q > 0, \quad \alpha > 0, \quad \epsilon > 0.
\end{equation*}
Una vez resuelta la LMI y obtenidos \((Q, Y)\), la ganancia se recupera como:
\begin{equation*}
    K = Y (Q)^{-1}.
\end{equation*}

Para comprobar la cota del error por elipsoide atractiva, partimos de:
\begin{equation*}
    \dot{V}(e) \leq -\alpha V(e) + \frac{1}{\epsilon}D^2.
\end{equation*}
Multiplicando por \(e^{\alpha t}\):
\begin{equation*}
    e^{\alpha t}\dot{V} + \alpha e^{\alpha t}V \leq \frac{1}{\epsilon}D^2 e^{\alpha t} \quad \Rightarrow \quad \frac{d}{dt}(e^{\alpha t}V) \leq \frac{1}{\epsilon}D^2 e^{\alpha t}.
\end{equation*}
Integrando de \(0\) a \(t\):
\begin{equation*}
    e^{\alpha t}V(t) - V(0) \leq \frac{D^2}{\epsilon \alpha}(e^{\alpha t} - 1),
\end{equation*}
\begin{equation*}
    V(t) \leq V(0)e^{-\alpha t} + \frac{D^2}{\epsilon \alpha}(1 - e^{-\alpha t}).
\end{equation*}
Tomando el límite cuando \(t \to \infty\):
\begin{equation*}
    \limsup_{t \to \infty} V(e(t)) \leq \frac{D^2}{\epsilon \alpha}.
\end{equation*}
Usando la desigualdad de Rayleigh:
\begin{equation*}
    \lambda_{\min}(P)\|e\|^2 \leq e^T P e = V(e) \quad \Rightarrow \quad \|e\| \leq \sqrt{\frac{V(e)}{\lambda_{\min}(P)}}.
\end{equation*}
Combinando con la cota superior de \(V\):
\begin{equation*}
    \|e\| \leq \sqrt{\frac{D^2}{\epsilon \alpha \lambda_{\min}(P)}}.
\end{equation*}
Por lo tanto, la elipsoide atractiva es:
\begin{equation*}
    \mathcal{E} = \left\{ e \in \mathbb{R}^{n+1} : e^T P e \leq \frac{D^2}{\epsilon \alpha} \right\}.
\end{equation*}





\end{itemize}









2.-  Dado el sistema

\begin{equation*}
    \begin{aligned}
        \dot{x}_1 =& cos(x_1+x_2+x_3)+x_2 \\
        \dot{x}_2 =& sin (x_1+x_2+x_3)+x_3\\
        \dot{x}_3=& sin(x_1x_2x_3+t)cos(x_1x_2x_3+t)+u\\
    y=& x_1    \end{aligned}
\end{equation*}


aplique un controlador por modos deslizantes mediante el método de elipsoides atractivas (SMC-AEM). (2 puntos)

\newpage







